% PTE Core 模考 1B —— 阅读 错题与详解·整理卷  版本 001（自包含）
% 编译: xelatex 本文件.tex   （需 Noto CJK + TeX Gyre Heros 字体）
% 本版本无外部图片，单文件即可复现。
\documentclass[11pt]{article}

% --------- 样式（原 ptestyle.sty，已内联）---------
% ============================================================
%  ptestyle.sty  —  PTE 模考错题整理卷 共享样式
%  编译引擎: XeLaTeX  (中英混排, 需要 Noto CJK 字体)
% ============================================================

\usepackage[a4paper,margin=2cm,headsep=4mm]{geometry}
\usepackage{fontspec}
\usepackage{xeCJK}
\usepackage[table]{xcolor}
\usepackage[normalem]{ulem}        % \sout 删除线
\usepackage{tabularx}
\usepackage{array}
\usepackage{ragged2e}
\usepackage{enumitem}
\usepackage[most]{tcolorbox}
\usepackage{fancyhdr}
\usepackage{graphicx}
\usepackage{pifont}                % \ding{51}=✓ \ding{55}=✗（阅读对错标注）
\usepackage{needspace}
\usepackage{tikz}

% ---------- 字体 ----------
\setmainfont{TeX Gyre Heros}            % 拉丁: Helvetica 风格(纤细清晰)
\setCJKmainfont{Noto Sans CJK SC}
\setCJKsansfont{Noto Sans CJK SC}
\newfontfamily\cjkserif{Noto Serif CJK SC}
\linespread{1.18}
\setlength{\parindent}{0pt}
\setlength{\parskip}{4pt}

% ---------- 配色 ----------
\definecolor{cWrong}{HTML}{D32F2F}     % 红  = 修改 / 错误
\definecolor{cFix}{HTML}{1B7F3B}       % 绿  = 正确建议
\definecolor{cDelete}{HTML}{1565C0}    % 蓝  = 删除
\definecolor{cInsert}{HTML}{7B1FA2}    % 紫  = 插入
\definecolor{cPartial}{HTML}{E07B00}   % 橙  = 部分得分
\definecolor{cNote}{HTML}{0277BD}      % 蓝  = 注解
\definecolor{cAccent}{HTML}{16A085}    % 主题青绿(平台主色)
\definecolor{cWrongBg}{HTML}{FBE0E0}
\definecolor{cDelBg}{HTML}{DCE8FB}
\definecolor{cInsBg}{HTML}{EEE0F5}
\definecolor{cHead}{HTML}{20303A}
\definecolor{cGrayBox}{HTML}{F4F6F7}
\definecolor{cModelBg}{HTML}{EAF7EF}
\definecolor{cYouBg}{HTML}{FFF7F5}
\definecolor{cNoteBg}{HTML}{EAF3FA}
\definecolor{cRule}{HTML}{D7DEE2}

% ---------- 行内批改宏 (复刻平台标注) ----------
% \fix{原词}{正确}  红色删除线 + 绿色上标改正  (修改/拼写/空格)
\newcommand{\fix}[2]{%
  \colorbox{cWrongBg}{\textcolor{cWrong}{\sout{#1}}}%
  \,\textsuperscript{\textcolor{cFix}{\footnotesize #2}}}
% \del{原词}  蓝色删除线  (应删除)
\newcommand{\del}[1]{%
  \colorbox{cDelBg}{\textcolor{cDelete}{\sout{#1}}}%
  \,\textsuperscript{\textcolor{cDelete}{\footnotesize 删}}}
% \ins{词}  紫色下划线  (应插入)
\newcommand{\ins}[1]{%
  \textsuperscript{\textcolor{cInsert}{\footnotesize $\wedge$}}%
  \colorbox{cInsBg}{\textcolor{cInsert}{#1}}}
% \miss{词}  灰色删除线  (口语漏读 / 未作答)
\newcommand{\miss}[1]{\textcolor{gray}{\sout{#1}}}
% \add{原词}{改正}  橙色波浪线 + "补"标签  (平台漏标、本卷补标的错误)
\definecolor{cAdd}{HTML}{E65100}
\newcommand{\add}[2]{%
  \textcolor{cAdd}{\uwave{#1}}%
  \,\textsuperscript{\textcolor{cAdd}{\footnotesize 补：#2}}}
% 高亮(口语发音): 好/中/差
\newcommand{\spk}[2]{\textcolor{#1}{#2}}

% ---------- 评分单元格着色 ----------
\newcommand{\sfull}[1]{\textcolor{cFix}{\bfseries #1}}     % 满分
\newcommand{\spart}[1]{\textcolor{cPartial}{\bfseries #1}} % 部分
\newcommand{\szero}[1]{\textcolor{cWrong}{\bfseries #1}}   % 零/低分

% ---------- 分数小标签 ----------
\newtcbox{\chip}[1][cAccent]{on line, boxsep=2pt, left=4pt, right=4pt, top=1pt,
  bottom=1pt, arc=3pt, colback=#1, colframe=#1, coltext=white, nobeforeafter,
  fontupper=\bfseries\footnotesize}

% ---------- 题块小标题 ----------
\newcommand{\ptesub}[1]{%
  \par\needspace{2\baselineskip}\smallskip
  {\color{cAccent}\rule[-0.2ex]{3pt}{1.05em}}\hspace{6pt}%
  {\bfseries\color{cHead}#1}\par\nopagebreak\smallskip}

% ---------- 题块外框 ----------
% \begin{qblock}{标号·题型}{右上信息(得分/用时)} ... \end{qblock}
\newtcolorbox{qblock}[2]{
  breakable, enhanced, arc=2.5pt,
  colback=white, colframe=cRule, boxrule=0.6pt,
  left=11pt, right=11pt, top=8pt, bottom=10pt,
  toptitle=3pt, bottomtitle=3pt, lefttitle=11pt, righttitle=11pt,
  colbacktitle=cHead, coltitle=white, fonttitle=\bfseries,
  title={#1\hfill\normalfont\bfseries\small #2}}

% ---------- 文本块: 题干 / 范文 / 你的作答 / 建议 ----------
\newtcolorbox{promptbox}{enhanced, breakable, colback=cGrayBox, colframe=cGrayBox,
  arc=2pt, left=8pt, right=8pt, top=6pt, bottom=6pt, boxrule=0pt}
\newtcolorbox{modelbox}{enhanced, breakable, colback=cModelBg, colframe=cModelBg,
  borderline west={3pt}{0pt}{cFix}, arc=1pt, left=8pt, right=8pt, top=6pt, bottom=6pt, boxrule=0pt}
\newtcolorbox{youbox}{enhanced, breakable, colback=cYouBg, colframe=cYouBg,
  borderline west={3pt}{0pt}{cWrong}, arc=1pt, left=8pt, right=8pt, top=6pt, bottom=6pt, boxrule=0pt}
\newtcolorbox{notebox}{enhanced, breakable, colback=cNoteBg, colframe=cNoteBg,
  borderline west={3pt}{0pt}{cNote}, arc=1pt, left=8pt, right=8pt, top=6pt, bottom=6pt, boxrule=0pt}

% ---------- 评分表 ----------
% 用法见正文; 这里给表头宏
\newcommand{\scorehead}{%
  \rowcolors{2}{cGrayBox}{white}
  \arrayrulecolor{cRule}}

% ---------- 章节大标题 ----------
\newcommand{\ptesection}[3]{% #1=中文名 #2=英文名 #3=该项分数
  \needspace{6\baselineskip}
  \begin{tcolorbox}[enhanced, colback=cAccent, colframe=cAccent, sharp corners,
     left=14pt, right=14pt, top=10pt, bottom=10pt, boxrule=0pt]
    {\fontsize{20}{24}\selectfont\bfseries\color{white} #1\ \ }%
    {\large\color{white!85} #2}\hfill
    {\large\color{white}本项得分\ \ }{\fontsize{22}{24}\selectfont\bfseries\color{white} #3}%
    \ {\color{white!85}/ 90}
  \end{tcolorbox}}

% ---------- 题型小节分隔 (口语 RA/RS/DI/RTS/ASQ 等) ----------
\newcommand{\ptetask}[2]{%
  \needspace{4\baselineskip}\medskip
  \begin{tcolorbox}[enhanced, colback=cHead, colframe=cHead, sharp corners,
     left=12pt, right=12pt, top=5pt, bottom=5pt, boxrule=0pt]
    {\large\bfseries\color{white} #1}\ \ {\color{white!80} #2}
  \end{tcolorbox}\smallskip}

% ---------- 口语 AI 语音识别着色 (复刻平台: 优/良/差 + 停顿) ----------
\definecolor{cSpkGood}{HTML}{2E9E5B}   % 优 绿
\definecolor{cSpkOk}{HTML}{C77F00}     % 良 黄/琥珀
\definecolor{cSpkBad}{HTML}{E0352B}    % 差 红
\newcommand{\sg}[1]{\textcolor{cSpkGood}{#1}}   % 优
\newcommand{\sy}[1]{\textcolor{cSpkOk}{#1}}     % 良
\newcommand{\sr}[1]{\textcolor{cSpkBad}{#1}}    % 差
\newcommand{\smiss}[1]{\textcolor{cSpkBad}{\sout{#1}}}  % 漏读 (红删除线)
\newcommand{\pse}{{\color{gray}\,/\,}}          % 停顿
\newcommand{\asrlegend}{{\footnotesize\color{gray}AI 语音识别\quad
  {\color{cSpkGood}\textbullet}\,优\ \ {\color{cSpkOk}\textbullet}\,良\ \
  {\color{cSpkBad}\textbullet}\,差\ \ {\color{gray}/}\,停顿\ \
  {\color{cSpkBad}\sout{词}}\,漏读}}
\newtcolorbox{asrbox}{enhanced, breakable, colback=white, colframe=cRule,
  boxrule=0.6pt, arc=1pt, left=8pt, right=8pt, top=6pt, bottom=6pt}

% ---------- 中文翻译行 ----------
\newcommand{\zh}[1]{{\footnotesize\textcolor{cNote}{\textbf{译}}\ \textcolor{gray}{#1}}}

% ---------- 阅读填空 / 选择 对错标注 ----------
\newcommand{\rok}[1]{{\color{cFix}\ding{51}\,#1}}                 % 选对(绿勾 + 词)
\newcommand{\rno}[2]{{\color{cWrong}\ding{55}\,\sout{#1}}\,{\footnotesize\color{cFix}(答案:\,#2)}}  % 选错: 你的(红删) + (答案:正确)
\newcommand{\rblank}[1]{{\color{cFix}\underline{\,#1\,}}}         % 直接给空的正确答案(无你的作答时)

% ---------- 颜色图例 ----------
\newcommand{\ptelegend}{%
  \begin{tcolorbox}[enhanced, colback=cGrayBox, colframe=cRule, boxrule=0.5pt,
    arc=2pt, left=10pt, right=10pt, top=6pt, bottom=6pt]
  {\bfseries\color{cHead}颜色说明\quad}
  \fix{改}{正}~修改/拼写\quad
  \del{删}~应删除\quad
  \ins{插}~应插入\quad
  \add{补}{改正}~平台漏标(本卷补标)\quad
  {\color{cFix}\rule{1em}{0.6em}}~范文/正确\quad
  {\color{cNote}\rule{1em}{0.6em}}~AI 建议
  \end{tcolorbox}}

% ---------- 页眉页脚 ----------
\pagestyle{fancy}
\fancyhf{}
\renewcommand{\headrulewidth}{0.4pt}
\renewcommand{\footrulewidth}{0pt}
\fancyhead[L]{\footnotesize\color{cHead}\bfseries PTE Core 模考 1B}
\fancyhead[R]{\footnotesize\color{gray} 错题与详解整理卷}
\fancyfoot[C]{\footnotesize\color{gray}\thepage}

\begin{document}

% --------- 封面（原 cover.tex）---------
% ---------- 封面 / 成绩总览 ----------
\definecolor{mListen}{HTML}{1F3A5F}
\definecolor{mRead}{HTML}{C2D70B}
\definecolor{mSpeak}{HTML}{4D4D4D}
\definecolor{mWrite}{HTML}{A01B7D}

\begin{titlepage}
\thispagestyle{empty}
\centering
\vspace*{1.4cm}

{\fontsize{30}{36}\selectfont\bfseries\color{cHead} PTE Core 模考 1B}\\[4mm]
{\fontsize{17}{20}\selectfont\color{cAccent}\bfseries 错题与详解 · 整理卷}\\[3mm]
{\color{gray} APEUni / 猩际 PTE 模考　·　提交时间 2026-06-21 12:36}\\[1.2cm]

% ---- 总分 ----
\begin{tcolorbox}[width=0.62\textwidth, halign=center, enhanced,
  colback=cAccent, colframe=cAccent, arc=6pt, boxrule=0pt, top=10pt, bottom=10pt]
  {\color{white}\large 总分 Overall}\\[2pt]
  {\fontsize{46}{50}\selectfont\bfseries\color{white} 40}
\end{tcolorbox}

\vspace{1.0cm}

% ---- 四项分数条形图 ----
\begin{tcolorbox}[width=0.84\textwidth, enhanced, colback=white, colframe=cRule,
  boxrule=0.8pt, arc=4pt, left=18pt, right=18pt, top=14pt, bottom=14pt]
\setlength{\tabcolsep}{6pt}\renewcommand{\arraystretch}{1.7}
\sffamily
\begin{tabular}{@{}r@{\hskip 8pt}l@{\hskip 10pt}l@{}}
 \bfseries Listening 听力 & {\color{mListen}\rule{3.67cm}{0.42cm}} & {\bfseries\large 33}\\
 \bfseries Reading 阅读   & {\color{mRead}\rule{4.67cm}{0.42cm}}   & {\bfseries\large 42}\\
 \bfseries Speaking 口语  & {\color{mSpeak}\rule{3.22cm}{0.42cm}}  & {\bfseries\large 29}\\
 \bfseries Writing 写作   & {\color{mWrite}\rule{6.33cm}{0.42cm}}  & {\bfseries\large 57}\\
\end{tabular}\\[2mm]
{\footnotesize\color{gray} 满分 90　|　刻度: 条长 = 得分 / 90}
\end{tcolorbox}

\vfill
\begin{flushleft}\small\color{gray}
\textbf{说明:}本卷由本次模考的题目与"详解"截图逐题誊写、重新排版而成,完整保留每道题的
\emph{题干 / 原文、范文、你的作答、逐项评分、彩色批改与 AI 中文建议}。\\
所有错误按平台标注用颜色标出(见下方图例)。原始"详解"PDF 不可复制、仅截图,本卷内容均为人工对照誊录,若与原图有出入以原图为准。
\end{flushleft}
\vspace{4mm}
\ptelegend
\end{titlepage}

% --------- 阅读正文（原 sec_reading_body.tex）---------
% ===================== 阅读 Reading =====================
\ptesection{阅读}{Reading · FIB + MCM + RO + MCS}{42}

\vspace{1mm}
{\small\color{gray} 本部分依据「阅读1 / 阅读2」逐题誊录。题型：\textbf{FIB} 读写填空（下拉选词）、\textbf{FIB(R)} 读填空（拖词入空）、
\textbf{MCM} 多选、\textbf{RO} 段落重排、\textbf{MCS} 单选。每题保留：\textbf{文章（含填空/选择答案）+ 中文翻译 + 选项 + 平台中文解析（根据原文 + 解题思路）+ 评分}。}
\vspace{1mm}

\begin{notebox}
{\footnotesize 标注：填空/选择处 {\color{cFix}\ding{51}\,绿=选对}；{\color{cWrong}\ding{55}\,红删除线=选错}，其后 {\footnotesize\color{cFix}(答案: 正确词)} 给正确答案。
选项里**加粗**的是正确项。共 16 题全部收录（阅读1 = Q1–Q10，阅读2 = Q11–Q16）。}
\end{notebox}

% =================== Q1 ===================
\begin{qblock}{Q1\quad FIB \#74\ \ Good Looks in Votes}{读写填空 · 评分 \textcolor{cAccent}{\bfseries 2 / 4}}
\ptesub{文章 Passage（含填空作答）}
\begin{promptbox}
It is tempting to try to prove that good looks win votes, and many academics have tried. The \rok{difficulty} is that beauty is in the eye of the \rno{audience}{beholder}, and you cannot behold a politician's face without a veil of extraneous prejudice getting in the way. Does George Bush possess a disarming grin, or a facetious \rno{binge}{smirk}? It's hard to find anyone who can look at the president without assessing him politically as well as \rok{physically}.
\end{promptbox}
\zh{想证明“颜值能赢得选票”很有诱惑力，许多学者也尝试过。难点在于“情人眼里出西施”——你看一位政客的脸时，总有一层无关的偏见挡在中间。乔治·布什是有一种令人放松的笑容，还是一种假笑？很难找到谁能在看总统时，不在政治层面、同时也在外貌层面对他做出评判。}
\ptesub{选项 Choices}
{\small
1.\ principle,\ idea,\ \textbf{difficulty},\ concept\\
2.\ people,\ \textbf{beholder},\ builder,\ audience\\
3.\ smell,\ complexion,\ \textbf{smirk},\ binge\\
4.\ culturally,\ \textbf{physically},\ economically,\ individually}
\ptesub{解析 Explanation}
\begin{notebox}\footnotesize
\textbf{1.\ difficulty} —— 根据原文：The \_\_\_ is that beauty is in the eye of the beholder…；此处需名词表“难处”，故 difficulty；principle 原则、idea 想法、concept 概念语意不合。\textbf{思路：单句理解}。\\[2pt]
\textbf{2.\ beholder} —— 固定搭配 \emph{beauty is in the eye of the beholder}（情人眼里出西施）；people 人们、builder 建造者、audience 观众均不合。\textbf{思路：固定搭配}。\\[2pt]
\textbf{3.\ smirk} —— 根据原文：Does George Bush possess a disarming grin, or a facetious \_\_\_? 与 grin 并列、需含贬义的名词，故 smirk（假笑）；smell 气味、complexion 肤色、binge 狂欢不合。\textbf{思路：单句理解}。\\[2pt]
\textbf{4.\ physically} —— 根据原文：…assessing him politically as well as \_\_\_；与 politically 并列，按逻辑选 physically；culturally 文化上、economically 经济上、individually 个人地不合。\textbf{思路：上下文逻辑}。
\end{notebox}
\end{qblock}

% =================== Q2 ===================
\begin{qblock}{Q2\quad FIB \#50\ \ Underground Houses}{读写填空 · 评分 \textcolor{cAccent}{\bfseries 2 / 5}}
\ptesub{文章 Passage（含填空作答）}
\begin{promptbox}
Underground houses have many advantages over conventional housing. Unlike conventional homes, they can be built on \rno{overhead}{steep} surfaces and can maximize space in small areas by going below the surface. In addition, the materials excavated in construction can be used in the building process. Underground houses have less surface area so fewer building materials are used, and \rok{maintenance} costs are lower. They are also wind, fire, and earthquake resistant, providing a secure and safe environment in extreme weather. One of the greatest benefits of underground living is energy \rno{consistency}{efficiency}. The earth's subsurface temperature remains stable, so underground dwellings benefit from geothermal mass and heat exchange, staying cool in the summer and warm in the winter. This saves around 80\% in energy costs. By \rno{inventing}{incorporating} solar design this energy bill \rok{can be reduced} to zero, providing hot water and heat to the home all year round.
\end{promptbox}
\zh{地下房屋相比传统住宅有许多优势。与传统住宅不同，它们可以建在陡峭的地面上，并通过向地下发展、在狭小空间里最大化利用面积。此外，施工中挖出的材料还能用于建造过程。地下房屋表面积更小，用的建筑材料更少，维护成本也更低。它们还能抗风、抗火、抗震，在极端天气下提供安全的环境。地下居住最大的好处之一是能源效率：地表以下温度稳定，地下住宅因而受益于地热与热交换，夏凉冬暖，可节省约 80\% 的能源开支。通过引入太阳能设计，这笔能源账单可被降至零，全年为家庭提供热水和供暖。}
\ptesub{选项 Choices}
{\small
1.\ geometric,\ flat,\ overhead,\ \textbf{steep}\\
2.\ \textbf{maintenance},\ heating,\ buoyancy,\ facility\\
3.\ ratio,\ consistency,\ \textbf{efficiency},\ renewal\\
4.\ intriguing,\ initiating,\ \textbf{incorporating},\ inventing\\
5.\ has reduced,\ \textbf{can be reduced},\ can reduce,\ has been reduced}
\ptesub{解析 Explanation}
\begin{notebox}\footnotesize
\textbf{1.\ steep} —— 根据原文：they can be built on \_\_\_ surfaces…；可建在“陡峭”的地面上，故 steep；geometric 几何的、flat 平的、overhead 头顶上的 不合。\textbf{思路：单句理解}。\\[2pt]
\textbf{2.\ maintenance} —— 根据原文：…fewer building materials are used, and \_\_\_ costs are lower；面积小、维修费用低，故 maintenance（维护）。\textbf{思路：单句理解}。\\[2pt]
\textbf{3.\ efficiency} —— 结合下文 \emph{This saves around 80\% in energy costs}（省约 80\% 能耗），故能源“效率”efficiency；ratio 比率、consistency 浓度、renewal 更新 不合。\textbf{思路：上下文理解}。\\[2pt]
\textbf{4.\ incorporating} —— 根据原文：By \_\_\_ solar design… providing hot water…；需“引入/纳入”太阳能设计，故 incorporating；intriguing 引发、initiating 发起、inventing 发明 不合。\textbf{思路：动宾搭配}。\\[2pt]
\textbf{5.\ can be reduced} —— 主语 energy bill 与“减少”是被动关系，且为假设、非已完成，故被动 can be reduced；has reduced、can reduce、has been reduced 不合。\textbf{思路：时态 + 被动语态}。
\end{notebox}
\end{qblock}

% =================== Q3 ===================
\begin{qblock}{Q3\quad FIB \#3\ \ Intelligence Comparison}{读写填空 · 评分 \textcolor{cAccent}{\bfseries 2 / 5}}
\ptesub{文章 Passage（含填空作答）}
\begin{promptbox}
Comparing the intelligence of animals of different species is difficult: how do you compare a dolphin and a horse? Psychologists have a technique for looking at intelligence that \rok{does} not require the cooperation of the animal involved. The relative size of an individual's brain is a reasonable indication of intelligence. Comparing \rno{with}{across} species is not as simple as generally expected. An elephant will have a larger brain than a human has simply because it is a large beast. \rno{Otherwise}{Instead}, we use the Cephalization index, which compares the size of an animal's brain with the size of its body. Based on the Cephalization index, the brightest animals on the planet are humans, \rok{followed} by great apes, porpoises and elephants. As a general \rno{principle}{rule}, animals that hunt for a living (like canines) are smarter than strict vegetarians (you don't need much intelligence to outsmart a leaf of lettuce). Animals that live in social groups are always smarter and have larger EQ's than solitary animals.
\end{promptbox}
\zh{比较不同物种动物的智力很难：你怎么比较一只海豚和一匹马？心理学家有一套观察智力的方法，它无需相关动物的配合——个体大脑的相对大小就是智力的一个合理指标。跨物种比较并不像一般以为的那么简单：大象的脑比人大，只是因为它体型大。于是我们改用“脑化指数”，把动物大脑的大小与其身体大小相比较。按脑化指数，地球上最聪明的动物是人类，其次是类人猿、海豚和大象。一般来说，以狩猎为生的动物（如犬科）比纯素食动物更聪明（要智胜一片生菜并不需要多少智力）。群居动物总是比独居动物更聪明、EQ 更高。}
\ptesub{选项 Choices}
{\small
1.\ can,\ do,\ did,\ \textbf{does}\\
2.\ \textbf{across},\ to,\ through,\ with\\
3.\ Then,\ \textbf{Instead},\ Because,\ Otherwise\\
4.\ \textbf{followed},\ follows,\ follow,\ following\\
5.\ theory,\ principal,\ \textbf{rule},\ principle}
\ptesub{解析 Explanation}
\begin{notebox}\footnotesize
\textbf{1.\ does} —— 根据原文：…a technique…that \_\_\_ not require the cooperation of the animal involved；需助动词、第三人称单数现在时，故 does；can/did/do 不匹配。\textbf{思路：主谓一致 + 时态}。\\[2pt]
\textbf{2.\ across} —— 根据原文：Comparing \_\_\_ species…；表“跨/不同物种间”比较，故 across；to/through/with 不合。\textbf{思路：单句理解}。\\[2pt]
\textbf{3.\ Instead} —— 根据原文：\_\_\_, we use the Cephalization index…；与前句形成转折，故副词 Instead；Then/Because/Otherwise 不合。\textbf{思路：单句理解（转折）}。\\[2pt]
\textbf{4.\ followed} —— 根据原文：…humans, \_\_\_ by great apes…；表“其次/紧随其后”，过去分词作后置定语，故 followed；follows/follow/following 不合。\textbf{思路：过去分词}。\\[2pt]
\textbf{5.\ rule} —— 固定搭配 \emph{as a general rule}（一般来说）；theory 理论、principal 校长、principle 法则 不合。\textbf{思路：固定搭配}。
\end{notebox}
\end{qblock}

% =================== Q4 ===================
\begin{qblock}{Q4\quad FIB \#70\ \ Mini Helicopter}{读写填空 · 评分 \textcolor{cAccent}{\bfseries 2 / 5}}
\ptesub{文章 Passage（含填空作答）}
\begin{promptbox}
A mini helicopter modelled on flying tree seeds could soon be flying overhead. Evan Ulrich and colleagues at the University of Maryland in College Park \rno{turned in}{turned to} the biological world for inspiration to build a scaled-down helicopter that could mimic the properties of full-size aircraft. The complex \rok{design} of full-size helicopters gets less efficient when shrunk, meaning that standard mini helicopters expend most of their power simply fighting to stay stable in the air. The researchers realized that a simpler aircraft designed to stay stable passively would use much less power and reduce manufacturing costs to boot. It turns out that nature \rno{is beating}{had beaten} them to it. The seeds of trees have a single-blade structure that \rok{allows} them to fly far away and drift safely to the ground. These seeds need no engine to \rno{fly}{spin} through the air, thanks to a process called autorotation. By analyzing the behavior of the samara with high-speed cameras, Ulrich and his team were able to copy its design.
\end{promptbox}
\zh{一种以飞行树种子为原型的微型直升机可能很快就会在头顶飞行。马里兰大学帕克分校的 Evan Ulrich 及同事转向生物界寻找灵感，想造一架能模仿全尺寸飞行器特性的缩小版直升机。全尺寸直升机复杂的设计在缩小后效率会下降，这意味着普通微型直升机的大部分动力都耗在勉力维持空中稳定上。研究者意识到，一架靠被动方式保持稳定的更简单飞行器会省电得多，还能降低制造成本。事实证明，大自然早已抢先一步：树的种子有一种单叶片结构，能让它们飞得很远并安全飘落到地面。借助一种叫“自转”的过程，这些种子无需发动机就能在空中旋转。通过用高速摄像机分析翅果（samara）的运动，Ulrich 团队得以复制其设计。}
\ptesub{选项 Choices}
{\small
1.\ \textbf{turned to},\ turned for,\ turned in,\ turned off\\
2.\ overhaul,\ gauge,\ imagination,\ \textbf{design}\\
3.\ is beating,\ was beaten,\ \textbf{had beaten},\ beaten\\
4.\ had allowed,\ allowed,\ \textbf{allows},\ will allow\\
5.\ \textbf{spin},\ fluctuate,\ drift,\ fly}
\ptesub{解析 Explanation}
\begin{notebox}\footnotesize
\textbf{1.\ turned to} —— 根据原文：…colleagues \_\_\_ the biological world for inspiration…；转向自然界寻求灵感，故 turned to（求助于）；turned in/for/off 不合。\textbf{思路：单句理解}。\\[2pt]
\textbf{2.\ design} —— 根据原文：The \_\_\_ of full-size helicopters gets less efficient when shrunk…；指“设计”，故 design；overhaul 检修、gauge 仪表、imagination 想象 不合。\textbf{思路：上下文理解}。\\[2pt]
\textbf{3.\ had beaten} —— 根据原文：It turns out that nature \_\_\_ them to it；\emph{beat sb. to it}（抢先一步），与 realized 等过去时呼应，表“过去的过去”，故过去完成 had beaten；is beating/was beaten/beaten 不合。\textbf{思路：时态}。\\[2pt]
\textbf{4.\ allows} —— 根据原文：…a single-blade structure that \_\_\_ them to fly far away…；定语从句谓语、一般现在时，故 allows；had allowed/allowed/will allow 时态不合。\textbf{思路：时态}。\\[2pt]
\textbf{5.\ spin} —— 根据原文：These seeds need no engine to \_\_\_ through the air, thanks to a process called autorotation；autorotation（自旋），故 spin（旋转）；fluctuate 波动、drift 漂流、fly 飞 不合。\textbf{思路：单句理解}。
\end{notebox}
\end{qblock}

% =================== Q5 ===================
\begin{qblock}{Q5\quad FIB \#290\ \ Power Mix}{读写填空 · 评分 \textcolor{cAccent}{\bfseries 3 / 5}}
\ptesub{文章 Passage（含填空作答）}
\begin{promptbox}
Imagine a time in the not too distant future when your power comes from a seamless mix of renewable energy and traditional sources. It is delivered by a grid that manages thousands of windmills and hundreds of thousands of customers. Computer \rno{controlling}{controlled}, the grid is able to manage instant variations in supply and demand and provides a real time power balance. Far more complex than anything \rok{in} existence today, it is the smart grid. This technology is a new frontier in power supply and seen as a green solution to current outdated management systems. When introduced, smart grids will result in energy savings and will allow consumers a choice in their electricity charges and to be able to select the cheapest time \rno{pins}{slots}. The difficulty for the energy industry is that smart grids \rok{do not exist} in reality and the power companies cannot experiment with existing supplies. Without an actual grid to conduct research on, Professor Wu has had to design a simulated laboratory including input from theoretical wind generators and solar panels to feed into a constantly operating system. For an authentic approach researchers built various types of equipment failures \rok{into} the grid to test the system. And it works.
\end{promptbox}
\zh{设想在不远的将来，你的电力来自可再生能源与传统能源的无缝组合，由一张管理着成千上万台风车、数十万用户的电网输送。在计算机控制下，电网能管理供需的瞬时变化，提供实时的电力平衡。它远比当今现存的任何系统都复杂——这就是智能电网。这项技术是电力供应的新前沿，被视为解决当前陈旧管理系统的绿色方案。智能电网一旦引入，将带来节能，并让消费者在电费上有选择权、能挑选最便宜的时段。能源行业的难处在于：智能电网在现实中并不存在，电力公司无法用现有供电做实验。由于没有真实电网可供研究，Wu 教授只得设计一个模拟实验室，引入理论风力发电机和太阳能板的输入，接入一个持续运行的系统。为贴近真实，研究者在电网中植入各种类型的设备故障来测试系统。结果——它奏效了。}
\ptesub{选项 Choices}
{\small
1.\ \textbf{controlled},\ has controlled,\ controls,\ controlling\\
2.\ with,\ without,\ of,\ \textbf{in}\\
3.\ cuts,\ pins,\ points,\ \textbf{slots}\\
4.\ does not exist,\ \textbf{do not exist},\ are not existing,\ not exist\\
5.\ \textbf{into},\ of,\ onto,\ above}
\ptesub{解析 Explanation}
\begin{notebox}\footnotesize
\textbf{1.\ controlled} —— 根据原文：Computer \_\_\_, the grid is able to…（独立主格结构）；Computer 与 control 为被动关系，用过去分词 controlled。\textbf{思路：独立主格 / 非谓语}。\\[2pt]
\textbf{2.\ in} —— 固定搭配 \emph{in existence}（现存的）：anything \_\_\_ existence today，故 in。\textbf{思路：介词 / 固定搭配}。\\[2pt]
\textbf{3.\ slots} —— 根据原文：…select the cheapest time \_\_\_；\emph{time slots}（时段），故 slots；cuts/pins/points 不合。\textbf{思路：词组搭配}。\\[2pt]
\textbf{4.\ do not exist} —— 根据原文：…smart grids \_\_\_ in reality…；主语复数、一般现在时否定，故 do not exist；does not exist/are not existing/not exist 不合。\textbf{思路：主谓一致}。\\[2pt]
\textbf{5.\ into} —— 根据原文：…built various types of equipment failures \_\_\_ the grid…；表“放入…其中”，故 into（与上文 feed into 呼应）；of/onto/above 不合。\textbf{思路：介词}。
\end{notebox}
\end{qblock}

% =================== Q6 ===================
\begin{qblock}{Q6\quad MCM-R \#52\ \ History of Sleep}{多选 Multiple Answers · 评分 \textcolor{cAccent}{\bfseries 0 / 3}}
\ptesub{文章 Passage}
\begin{promptbox}
September 2, 1752, was a great day in the history of sleep. That Wednesday evening, millions of British subjects in England and the colonies went peacefully to sleep and did not wake up until twelve days later. Behind this feat of narcoleptic prowess was not some revolutionary hypnotic technique or miraculous pharmaceutical discovered in the West Indies. It was, rather, the British Calendar Act of 1751, which declared the day after Wednesday 2nd to be Thursday 14th. Prior to that cataclysmic September evening, the official British calendar differed from that of continental Europe by eleven days---that is, September 2 in London was September 13 in Paris, Lisbon, and Berlin. The discrepancy had sprung from Britain's continued use of the Julian calendar, which had been the official calendar of Europe from its invention by Julius Caesar (after whom it was named) in 45 B.C. until the decree of Pope Gregory XIII in 1582. Caesar's calendar, which consisted of eleven months of 30 or 31 days and a 28-day February (extended to 29 days every fourth year), was actually quite accurate: it erred from the real solar calendar by only 11.5 minutes a year. After centuries, though, even a small inaccuracy like this adds up. By the sixteenth century, it had put the Julian calendar behind the solar one by 10 days. In Europe, in 1582, Pope Gregory XIII ordered the advancement of the Julian calendar by 10 days and introduced a new corrective device to curb further error: century years such as 1700 or 1800 would no longer be counted as leap years, unless they were (like 1600 or 2000) divisible by 400.
\end{promptbox}
\zh{1752 年 9 月 2 日是睡眠史上“伟大”的一天。那个周三晚上，英格兰及其殖民地的数百万英国臣民安然入睡，直到十二天后才醒来。这桩“嗜睡壮举”的背后，并非什么革命性的催眠术或在西印度群岛发现的神奇药物，而是 1751 年的《英国历法法案》——它宣布周三 2 日的次日即为周四 14 日。在那个翻天覆地的九月夜晚之前，英国官方历法与欧洲大陆相差十一天——伦敦的 9 月 2 日，就是巴黎、里斯本和柏林的 9 月 13 日。这一差异源于英国一直沿用儒略历：自公元前 45 年尤利乌斯·凯撒发明（历法因他得名）起，到 1582 年教皇格里高利十三世颁布敕令为止，它都是欧洲的官方历法。凯撒的历法由 11 个 30 或 31 天的月份和一个 28 天的二月（每四年延长为 29 天）组成，其实相当准确：每年与真实太阳历仅差 11.5 分钟。但几个世纪累积，即便这点小误差也会叠加；到 16 世纪，儒略历已比太阳历落后 10 天。1582 年，教皇下令把儒略历提前 10 天，并引入新纠错机制：1700、1800 这样的世纪年不再算闰年，除非（像 1600、2000）能被 400 整除。}
\ptesub{题目 Question}
\begin{notebox}\small What factors were involved in the disparity between the calendars of Britain and Europe in the 17th century?\end{notebox}
\ptesub{选项 Choices（多选）}
{\small
A）the provisions of the British Calendar Act of 1751\\
{\color{cFix}\textbf{B）Britain's continued use of the Julian calendar}}\\
{\color{cFix}\textbf{C）the accrual of very minor differences between the calendar used in Britain and real solar events}}\\
D）the failure to include years divisible by four as leap years\\
{\color{cFix}\textbf{E）the decree of Pope Gregory XIII}}\\
F）revolutionary ideas which had emerged from the West Indies\\
G）Britain's use of a calendar consisting of twelve months rather than eleven}
\smallskip\par
{\small 正确答案 \textcolor{cFix}{\bfseries B, C, E}\qquad 你的答案 {\color{cWrong}\bfseries C, F}\;{\footnotesize\color{gray}（C 对、F 错；多选错选不得分）}}
\end{qblock}

% =================== Q7 ===================
\begin{qblock}{Q7\quad MCM-R \#53\ \ X-ray}{多选 Multiple Answers · 评分 \textcolor{cAccent}{\bfseries 1 / 2}}
\ptesub{文章 Passage}
\begin{promptbox}
X-ray crystallography is the study of crystal structures through X-ray diffraction techniques. When an X-ray beam bombards a crystalline lattice in a given orientation, the beam is scattered in a definite manner characterized by the atomic structure of the lattice. This phenomenon, known as X-ray diffraction, occurs when the wavelength of X-rays and the interatomic distances in the lattice have the same order of magnitude. In 1912, the German scientist Max von Laue predicted that crystals exhibit diffraction qualities. Concurrently, W. Friedrich and P. Knipping created the first photographic diffraction patterns. A year later, Lawrence Bragg successfully analyzed the crystalline structures of potassium chloride and sodium chloride using X-ray crystallography, and developed a rudimentary treatment for X-ray/crystal interaction (Bragg's Law). Bragg's research provided a method to determine a number of simple crystal structures for the next 50 years. In the 1960s, the capabilities of X-ray crystallography were greatly improved by the incorporation of computer technology. Modern X-ray crystallography provides the most powerful and accurate method for determining single-crystal structures. Structures containing 100--200 atoms now can be analyzed on the order of 1--2 days, whereas before the 1960s a 20-atom structure required 1--2 years for analysis. Through X-ray crystallography the chemical structure of thousands of organic, inorganic, organometallic, and biological compounds are determined every year.
\end{promptbox}
\zh{X 射线晶体学是通过 X 射线衍射技术研究晶体结构的学科。当 X 射线束以特定取向轰击晶格时，射线会以由晶格原子结构决定的特定方式被散射；这种现象称为 X 射线衍射，当 X 射线波长与晶格中原子间距处于同一数量级时即发生。1912 年，德国科学家马克斯·冯·劳厄预言晶体会表现出衍射特性；同期 W. 弗里德里希与 P. 克尼平制作出首批衍射照片图样。一年后，劳伦斯·布拉格用 X 射线晶体学成功分析了氯化钾、氯化钠的晶体结构，并提出 X 射线/晶体相互作用的基本处理方法（布拉格定律）。布拉格的研究为之后 50 年测定若干简单晶体结构提供了方法。20 世纪 60 年代，计算机技术的引入大大提升了其能力。现代 X 射线晶体学提供了测定单晶结构最强大、最精确的方法：含 100–200 个原子的结构如今 1–2 天即可分析，而 60 年代前一个 20 原子的结构需 1–2 年。借助 X 射线晶体学，每年都能测定数以千计的有机、无机、有机金属和生物化合物的化学结构。}
\ptesub{题目 Question}
\begin{notebox}\small Which of the following factors are consistent with the theory of X-ray crystallography?\end{notebox}
\ptesub{选项 Choices（多选）}
{\small
A）X-ray crystallization causes a reduction in the interatomic distances of wavelengths.\\
{\color{cFix}\textbf{B）X-rays are scattered according to the atomic structure of the crystal lattice.}}\\
C）The analysis of chemical compounds was only possible after the development of computer technology.\\
{\color{cFix}\textbf{D）The process can be used to determine the chemical structure of biological compounds.}}\\
E）X-rays will not diffract in crystalline substances.}
\smallskip\par
{\small 正确答案 \textcolor{cFix}{\bfseries B, D}\qquad 你的答案 {\color{cWrong}\bfseries D}\;{\footnotesize\color{gray}（漏选 B）}}
\end{qblock}

% =================== Q8 ===================
\begin{qblock}{Q8\quad RO \#307\ \ Noise and Study}{段落重排 Re-order · 评分 \textcolor{cAccent}{\bfseries 1 / 3}}
\ptesub{段落 Paragraphs（按原乱序编号）}
\begin{promptbox}\small
\textbf{1)} However, one general rule for all students is that the television seems to be more of a distraction than music or other background noise, so leave the TV off when you are reading or studying. Also, don't let yourself distracted by computer games, email, or internet surfing.\\[2pt]
\textbf{2)} Some students say that they need complete quiet to read and study.\\[2pt]
\textbf{3)} The point is, you should know the level of noise that is optimal for your own studying.\\[2pt]
\textbf{4)} Others study best in crowded, noisy rooms because the noise actually helps them concentrate.
\end{promptbox}
\zh{1）不过，对所有学生都适用的一条通则是：电视似乎比音乐或其他背景噪音更让人分心，所以读书学习时把电视关掉，也别被电脑游戏、邮件或上网分心。2）有些学生说，他们读书学习需要绝对安静。3）关键在于，你要知道最适合自己学习的噪音水平。4）另一些人在拥挤吵闹的房间里学得最好，因为噪音反而帮他们集中注意力。}
\smallskip\par
{\small 正确顺序 \textcolor{cFix}{\bfseries 2 → 4 → 3 → 1}\qquad 你的顺序 {\color{cWrong}\bfseries 2 → 1 → 4 → 3}}
\end{qblock}

% =================== Q9 ===================
\begin{qblock}{Q9\quad RO \#300\ \ GPS Tracking（GPS 定位）}{段落重排 Re-order · 评分 \textcolor{cAccent}{\bfseries 3 / 3}}
\ptesub{段落 Paragraphs（按原乱序编号）}
\begin{promptbox}\small
\textbf{1)} The collars transmitted each animal's position every four hours, for up to two years.\\[2pt]
\textbf{2)} But in 2010 and 2011, Vanessa Hull of Michigan State University and her colleagues were given permission to attach GPS tracking collars to five pandas in the Wolong National Nature Reserve in China.\\[2pt]
\textbf{3)} We know very little about wild pandas because they are so rare and live in almost impenetrable forest.\\[2pt]
\textbf{4)} The team found that the home ranges of individual pandas overlapped and on a few occasions, two animals spent several weeks in close proximity.
\end{promptbox}
\zh{1）项圈每隔四小时传送一次每只动物的位置，持续长达两年。2）但在 2010 和 2011 年，密歇根州立大学的 Vanessa Hull 及同事获准给卧龙国家级自然保护区的五只大熊猫戴上 GPS 追踪项圈。3）我们对野生大熊猫知之甚少，因为它们非常稀少、且生活在几乎无法穿越的森林里。4）团队发现，各只熊猫的活动范围相互重叠，偶尔有两只熊猫连续数周近距离共处。}
\smallskip\par
{\small 正确顺序 \textcolor{cFix}{\bfseries 3 → 2 → 1 → 4}\qquad 你的顺序 \textcolor{cFix}{\bfseries 3 → 2 → 1 → 4}\;{\footnotesize\color{gray}（全对）}}
\par{\footnotesize\color{gray}思路：先“我们知之甚少（3）”，再“2010–11 装 GPS 项圈（2）”，接着“项圈传送位置（1）”，最后“发现活动范围重叠（4）”。}
\end{qblock}

% =================== Q10 ===================
\begin{qblock}{Q10\quad FIB(R) \#483\ \ Martens' Diet}{读填空（拖词）· 评分 \textcolor{cAccent}{\bfseries 2 / 4}}
\ptesub{文章 Passage（含填空作答）}
\begin{promptbox}
Studies of pine martens in Scotland have shown that the diet varies \rok{seasonally} with small mammals, berries (in late summer/autumn) and small birds being the main foods. Recent work on a plantation has shown that martens \rok{establish} their home ranges in areas dominated by forests and \rno{lately}{dense} shrubs. Within home ranges, martens utilize areas of grassy vegetation within the forest which are typically associated with Microtus voles, for which a strong selective \rno{decision}{preference} over other small mammals is shown.
\end{promptbox}
\zh{在苏格兰对松貂的研究表明，其食物随季节变化，以小型哺乳动物、浆果（夏末/秋季）和小鸟为主要食物。一处人工林的最新研究显示，松貂把活动范围建立在以森林和茂密灌木为主的区域。在活动范围内，松貂会利用森林中通常与田鼠（Microtus voles）相关的草地植被区域——相对其他小型哺乳动物，它们对田鼠表现出强烈偏好。}
\ptesub{选项 Choices（拖词词库，多于空格数）}
{\small \textbf{seasonally}\quad \textbf{establish}\quad \textbf{dense}\quad \textbf{preference}\quad ｜\quad lately,\ decision,\ complicated（多余干扰词）}
\ptesub{解析 Explanation}
\begin{notebox}\footnotesize
\textbf{1.\ seasonally} —— 根据原文：the diet varies \_\_\_ with small mammals…；副词修饰 varies，故 seasonally（季节性地）。\textbf{思路：词性 + 单句理解}。\\[2pt]
\textbf{2.\ establish} —— 根据原文：martens \_\_\_ their home ranges…；需动词谓语“建立”，故 establish。\textbf{思路：动词谓语}。\\[2pt]
\textbf{3.\ dense} —— 根据原文：forests and \_\_\_ shrubs；需形容词修饰 shrubs，故 dense（茂密的）；lately（副词）不合。\textbf{思路：词性 + 搭配}。\\[2pt]
\textbf{4.\ preference} —— 根据原文：a strong selective \_\_\_ over other small mammals…；固定搭配 \emph{a preference for A over B}（相比 B 更偏好 A），故 preference；decision 不合。\textbf{思路：固定搭配}。
\end{notebox}
\end{qblock}

% =================== Q11 ===================
\begin{qblock}{Q11\quad FIB(R) \#165\ \ Bees' Die-off}{读填空（拖词）· 评分 \textcolor{cAccent}{\bfseries 4 / 5}}
\ptesub{文章 Passage（含填空作答）}
\begin{promptbox}
It sounds like something out of a science fiction movie --- or nightmare: millions of honeybees \rok{suddenly} dying off, their bodies never found. Scientists have \rok{named} the phenomenon `Colony Collapse Disorder', but they aren't \rno{deliberating}{united} on the reason. Theories abound as to the \rok{cause} of the mass die-off, ranging from the unlikely (cellphones affecting bees' navigational abilities) to the more \rok{plausible} though still debated (widespread pesticide use).
\end{promptbox}
\zh{这听起来像科幻电影——或噩梦：数百万只蜜蜂突然死去，却找不到它们的尸体。科学家把这一现象命名为“蜂群崩溃失调症”，但对其原因尚未达成一致。关于这场大规模死亡的原因众说纷纭，从不太可能的（手机影响蜜蜂导航能力）到更可信、但仍有争议的（农药的广泛使用）。}
\ptesub{选项 Choices（拖词词库，多于空格数）}
{\small \textbf{suddenly}\quad \textbf{named}\quad \textbf{united}\quad \textbf{cause}\quad \textbf{plausible}\quad ｜\quad deliberating,\ possibility,\ authored（多余干扰词）}
\ptesub{解析 Explanation}
\begin{notebox}\footnotesize
\textbf{1.\ suddenly} —— …millions of honeybees \_\_\_ dying off…；副词修饰 dying，故 suddenly（突然）。\textbf{思路：词性}。\\[2pt]
\textbf{2.\ named} —— Scientists have \_\_\_ the phenomenon…；\emph{name sth}（命名），现在完成 have named。\textbf{思路：动词 + 时态}。\\[2pt]
\textbf{3.\ united} —— they aren't \_\_\_ on the reason；\emph{be united on}（在…上达成一致），故 united；deliberating（审议）不合。\textbf{思路：搭配}。\\[2pt]
\textbf{4.\ cause} —— as to the \_\_\_ of the mass die-off…；名词“原因”，故 cause。\textbf{思路：词性 + 单句理解}。\\[2pt]
\textbf{5.\ plausible} —— to the more \_\_\_ though still debated…；形容词“可信的”，故 plausible。\textbf{思路：词性 + 上下文}。
\end{notebox}
\end{qblock}

% =================== Q12 ===================
\begin{qblock}{Q12\quad FIB(R) \#481\ \ Contagious Emotions}{读填空（拖词）· 评分 \textcolor{cAccent}{\bfseries 3 / 3}}
\ptesub{文章 Passage（含填空作答）}
\begin{promptbox}
As research has shown, emotions are contagious. And empaths are especially \rok{sensitive} to others' emotional energies. Because they can get easily exhausted in crowds, be drawn into codependent \rok{relationships}, exhaust themselves trying to solve others' problems, or burn out from too much caregiving. Yet empathy is also a gift that brings greater \rok{insight} and understanding. Some of the finest therapists, doctors, nurses, professors, writers, designers, musicians, artists and leaders in many have been empaths.
\end{promptbox}
\zh{研究表明，情绪会传染。共情者（empaths）对他人的情感能量尤其敏感。他们在人群中容易精疲力竭，会陷入相互依赖的关系，会因试图解决别人的问题而耗尽自己，或因过度照顾他人而倦怠。然而，共情也是一种天赋，能带来更深的洞察与理解。许多最优秀的治疗师、医生、护士、教授、作家、设计师、音乐家、艺术家与领袖，都曾是共情者。}
\ptesub{选项 Choices（拖词词库，多于空格数）}
{\small \textbf{sensitive}\quad \textbf{relationships}\quad \textbf{insight}\quad ｜\quad confusion,\ issues,\ resistant（多余干扰词）}
\ptesub{解析 Explanation}
\begin{notebox}\footnotesize
\textbf{1.\ sensitive} —— …empaths are especially \_\_\_ to others' emotional energies；be 后接形容词，故 sensitive（敏感的）；resistant 不合。\textbf{思路：词性}。\\[2pt]
\textbf{2.\ relationships} —— …be drawn into codependent \_\_\_；形容词后需名词，故 relationships（关系）；issues 不合。\textbf{思路：词性}。\\[2pt]
\textbf{3.\ insight} —— …brings greater \_\_\_；greater 后需名词，故 insight（洞察）；confusion 不合。\textbf{思路：词性 + 上下文}。
\end{notebox}
\end{qblock}

% =================== Q13 ===================
\begin{qblock}{Q13\quad FIB(R) \#465\ \ Utopias}{读填空（拖词）· 评分 \textcolor{cAccent}{\bfseries 2 / 6}}
\ptesub{文章 Passage（含填空作答）}
\begin{promptbox}
Many Utopias have been dreamed up through the ages. From Plato's Republic to Thomas More's Utopia and beyond, serious thinkers have \rno{exchanged}{envisioned} societies where people live in peace and harmony. Most of these imaginary worlds have things in common: everybody is equal and plays a part in the running of the society; nobody goes without the essentials of life; people live mostly off the land; often there is no money, and so on. Another thing they have in \rok{common} is that, to the average person, they appear distasteful or unworkable since they do not take into account ordinary human nature or feelings. Architects have got in on the act, too. After the Great Fire of London, Christopher Wren drew up plans for a reconstruction of the whole city, including \rno{horizontal}{precise} grand avenues; in the 20th century there was Le Corbusier's Radiant City in which, if you weren't in a car or didn't have one, life would have been a nightmare. Also in the 20th century, another famous architect, Frank Lloyd Wright, \rok{dreamed} up a perfect city that got no further than the drawing-board. Wright believed that what was wrong with modern cities was, in his words, rent: ideas, land, even money itself, had to be paid for. He saw this as a form of slavery and believed that modern city dwellers had no sense of themselves as productive individuals. Thus, Wright's city was to be made up of numerous individual homesteads, and the houses themselves were to be simple, functional and in \rno{pieced}{harmony} with the environment. Everyone would own enough land to grow food for himself and his family. No outsiders would be allowed to come between the citizen and what he produced, or to exploit both for money. Goods and services would all be \rno{envisioned}{exchanged}, not bought and sold for profit.
\end{promptbox}
\zh{古往今来人们构想过许多乌托邦。从柏拉图《理想国》到托马斯·莫尔《乌托邦》乃至之后，严肃的思想家们设想过种种人们和平和谐共处的社会。这些想象世界大多有共同点：人人平等、都参与社会运作；无人缺乏生活必需品；人们多半靠土地为生；往往没有货币，等等。它们还有一个共同点：在普通人看来令人反感或行不通，因为没有考虑普通人的本性与情感。建筑师们也参与其中。伦敦大火后，克里斯托弗·雷恩制定了重建全城的规划，包括精确规整的大道；20 世纪有勒·柯布西耶的“光辉城市”，在那里若你没有汽车，生活会是噩梦。同样在 20 世纪，另一位著名建筑师弗兰克·劳埃德·赖特构想了一座完美城市，却止步于图纸。赖特认为现代城市的问题在于——用他的话说——租金：理念、土地乃至金钱本身都要花钱买。他视之为一种奴役，认为现代城市居民没把自己当作有生产力的个体。因此，赖特的城市将由许多个人家园组成，房屋本身简单、实用、与环境和谐。人人拥有足够土地为自己和家人种粮；不允许外人介入公民与其产出之间、或为牟利剥削两者。商品和服务都将被交换，而非为利润买卖。}
\ptesub{选项 Choices（拖词词库，多于空格数）}
{\small \textbf{envisioned}\quad \textbf{common}\quad \textbf{precise}\quad \textbf{dreamed}\quad \textbf{harmony}\quad \textbf{exchanged}\quad ｜\quad regulated,\ paved,\ incorporated（多余干扰词）}
\ptesub{解析 Explanation}
\begin{notebox}\footnotesize
\textbf{1.\ envisioned} —— serious thinkers have \_\_\_ societies…；动词谓语“构想/设想”，故 envisioned；exchanged（交换）不合。\textbf{思路：动词 + 上下文}。\\[2pt]
\textbf{2.\ common} —— Another thing they have in \_\_\_…；\emph{have in common}（有共同点），故 common。\textbf{思路：固定搭配}。\\[2pt]
\textbf{3.\ precise} —— including \_\_\_ grand avenues；形容词“精确规整的”，故 precise；horizontal（水平的）不合。\textbf{思路：词性 + 语义}。\\[2pt]
\textbf{4.\ dreamed} —— Wright \_\_\_ up a perfect city；\emph{dream up}（凭空设想），故 dreamed。\textbf{思路：动词短语}。\\[2pt]
\textbf{5.\ harmony} —— in \_\_\_ with the environment；\emph{in harmony with}（与…和谐），故 harmony；pieced 不合。\textbf{思路：固定搭配}。\\[2pt]
\textbf{6.\ exchanged} —— Goods and services would all be \_\_\_…；被动语态、“交换”，故 exchanged；envisioned 不合。\textbf{思路：被动语态 + 上下文}。
\end{notebox}
\end{qblock}

% =================== Q14 ===================
\begin{qblock}{Q14\quad FIB(R) \#449\ \ Supply and Demand}{读填空（拖词）· 评分 \textcolor{cAccent}{\bfseries 3 / 4}}
\ptesub{文章 Passage（含填空作答）}
\begin{promptbox}
The supply of a thing, in the phrase `supply and demand', is the amount that will be offered for sale at each of a series of prices; the demand is the amount that will be bought at each of a series of prices. The principle that value depends on supply and demand means that in the case of nearly every commodity, more will be bought if the price is lowered; less will be bought if the price is \rok{raised}. Therefore sellers, if they wish to induce buyers to take more of a commodity than they are already doing, must reduce its price; if they raise its price, they will sell less. If there is a general falling-off in demand --- due, say, to trade depression --- sellers will either have to \rok{reduce} prices or put less on the \rok{market}; they will not be able to sell the same \rno{gear}{amount} at the same price. Similarly with supply. At a certain price a certain amount will be offered for sale; at a higher price more will be offered; at a lower price less. If consumers want more, they must offer a higher price; if they want less, they will probably be able to force prices down. That is the first result of a change in demand or supply.
\end{promptbox}
\zh{在“供给与需求”这一短语中，某物的供给是指在一系列价格下愿意出售的数量；需求是指在一系列价格下会被买进的数量。“价值取决于供求”这一原理意味着：对几乎每种商品而言，价格降低则买得更多，价格提高则买得更少。因此，卖方若想让买家比现在买得更多，就必须降价；若提价，就会卖得更少。若需求普遍下滑——比如因贸易萧条——卖方要么降价，要么减少投放市场的量；他们将无法以同样的价格卖出同样的数量。供给方面也类似：某一价格下会有一定数量供售，价格越高供应越多，越低越少。消费者想要更多就得出更高的价；想要更少则多半能把价格压下来。这就是供求变化的首要结果。}
\ptesub{选项 Choices（拖词词库，多于空格数）}
{\small \textbf{raised}\quad \textbf{reduce}\quad \textbf{market}\quad \textbf{amount}\quad ｜\quad admit,\ recorded,\ rate,\ gear（多余干扰词）}
\ptesub{解析 Explanation}
\begin{notebox}\footnotesize
\textbf{1.\ raised} —— less will be bought if the price is \_\_\_；与 lowered 相对，故 raised（提高）。\textbf{思路：反义对应}。\\[2pt]
\textbf{2.\ reduce} —— sellers will either have to \_\_\_ prices…；动词“降低”，故 reduce。\textbf{思路：动词}。\\[2pt]
\textbf{3.\ market} —— put less on the \_\_\_；\emph{put on the market}（投放市场），故 market。\textbf{思路：固定搭配}。\\[2pt]
\textbf{4.\ amount} —— sell the same \_\_\_ at the same price；same 后需名词“数量”，且与 sell 搭配，故 amount；gear（齿轮）不合。\textbf{思路：动宾搭配}。
\end{notebox}
\end{qblock}

% =================== Q15 ===================
\begin{qblock}{Q15\quad MCS-R \#118\ \ Writing in College}{单选 Single Answer · 评分 \textcolor{cAccent}{\bfseries 0 / 1}}
\ptesub{文章 Passage}
\begin{promptbox}
One of the first things you'll discover as a college student is that writing in college is different from writing in high school. Certainly a lot of what your high school writing teachers taught you will be useful to you as you approach writing in college: you will want to write clearly, to have an interesting and arguable thesis, to construct paragraphs that are coherent and focused, and so on. Still, many students enter college relying on writing strategies that served them well in high school but that won't serve them well here. Old formulae, such as the five-paragraph theme, aren't sophisticated or flexible enough to provide a sound structure for a college paper. And many of the old tricks --- such as using elevated language, or repeating yourself so that you might meet a ten-page requirement --- will fail you now.
\end{promptbox}
\zh{你上大学后最先发现的事情之一，是大学写作与高中写作不同。诚然，高中写作老师教你的很多东西在大学写作中仍然有用：你会想把文章写清楚、提出有趣且可论辩的论点、写出连贯而聚焦的段落，等等。然而，许多学生带着在高中行之有效、却在这里行不通的写作策略进入大学。诸如“五段式”这样的老套路，不够精巧灵活，无法为大学论文提供合理结构。许多老把戏——比如堆砌华丽辞藻，或为凑满十页要求而反复啰嗦——如今都会让你失败。}
\ptesub{题目 Question}
\begin{notebox}\small According to the writer, a student might repeat himself to \_\_\_\_.\end{notebox}
\ptesub{选项 Choices（单选）}
{\small
A）write a conclusion for the essay\\
B）remind the teacher of what he has written\\
{\color{cFix}\textbf{C）increase the length of essay}}\\
D）emphasize the main argument of the essay}
\smallskip\par
{\small 正确答案 \textcolor{cFix}{\bfseries C}\qquad 你的答案 {\color{cWrong}\bfseries D}\;{\footnotesize\color{gray}（文中 repeating yourself so that you might meet a ten-page requirement——重复是为了凑页数，即增加篇幅，故 C）}}
\end{qblock}

% =================== Q16 ===================
\begin{qblock}{Q16\quad MCS-R \#14\ \ Art}{单选 Single Answer · 评分 \textcolor{cAccent}{\bfseries 0 / 1}}
\ptesub{文章 Passage}
\begin{promptbox}
Many argue that art cannot be defined. We could go about this in several ways. Art is often considered as the process or product of deliberately arranging elements in a way that appeals to the senses or emotions. It encompasses a diverse range of human activities, creations and ways of expression, including music, literature, film, sculpture and paintings. The meaning of art is explored in a branch of philosophy known as aesthetics. At least, that is what Wikipedia claims.
\end{promptbox}
\zh{许多人认为艺术无法被定义。我们可以从几个角度来看。艺术常被视为有意排列各种元素、以打动感官或情感的过程或产物。它涵盖多种多样的人类活动、创作与表达方式，包括音乐、文学、电影、雕塑和绘画。艺术的含义在一个被称为“美学”的哲学分支中被探讨。至少，维基百科是这么说的。}
\ptesub{题目 Question}
\begin{notebox}\small What is the main idea of the passage?\end{notebox}
\ptesub{选项 Choices（单选）}
{\small
{\color{cFix}\textbf{A）Art is a difficult and complex form to explain.}}\\
B）Wikipedia defines art under aesthetics, which is a branch of philosophy.\\
C）Music, literature, film and sculpture do not define art.\\
D）Art is directed in a way that it deliberately appeals to the emotions of people.}
\smallskip\par
{\small 正确答案 \textcolor{cFix}{\bfseries A}\qquad 你的答案 {\color{cWrong}\bfseries B}\;{\footnotesize\color{gray}（全文围绕“艺术难以定义”，主旨为 A；B/C/D 为细节或片面）}}
\end{qblock}

\end{document}
